\documentclass[10pt,a4paper]{amsart}
\usepackage{amsmath,amssymb,amsthm,graphicx}
% For commutative diagrams we use Xy-pic
\usepackage[all]{xy}

% We use 2cell for 2-commutative diagrams.
\xyoption{2cell}
% Theorem environments.
%
\theoremstyle{plain}
\newtheorem{theorem}[subsection]{Theorem}
\newtheorem{proposition}[subsection]{Proposition}
\newtheorem{lemma}[subsection]{Lemma}

\theoremstyle{definition}
\newtheorem{definition}[subsection]{Definition}
\newtheorem{example}[subsection]{Example}
\newtheorem{exercise}[subsection]{Exercise}
\newtheorem{situation}[subsection]{Situation}

\theoremstyle{remark}
\newtheorem{remark}[subsection]{Remark}
\newtheorem{remarks}[subsection]{Remarks}

\numberwithin{equation}{subsection}

% Macros
%
\def\lim{\mathop{\mathrm{lim}}\nolimits}
\def\colim{\mathop{\mathrm{colim}}\nolimits}
\def\Spec{\mathop{\mathrm{Spec}}}
\def\Hom{\mathop{\mathrm{Hom}}\nolimits}
\def\Ext{\mathop{\mathrm{Ext}}\nolimits}
\def\SheafHom{\mathop{\mathcal{H}\!\mathit{om}}\nolimits}
\def\SheafExt{\mathop{\mathcal{E}\!\mathit{xt}}\nolimits}
\def\Sch{\mathit{Sch}}
\def\Mor{\mathop{\mathrm{Mor}}\nolimits}
\def\Ob{\mathop{\mathrm{Ob}}\nolimits}
\def\Sh{\mathop{\mathit{Sh}}\nolimits}
\def\NL{\mathop{N\!L}\nolimits}
\def\CH{\mathop{\mathrm{CH}}\nolimits}
\def\proetale{{pro\text{-}\acute{e}tale}}
\def\etale{{\acute{e}tale}}
\def\QCoh{\mathit{QCoh}}
\def\Ker{\mathop{\mathrm{Ker}}}
\def\Im{\mathop{\mathrm{Im}}}
\def\Coker{\mathop{\mathrm{Coker}}}
\def\Coim{\mathop{\mathrm{Coim}}}

% Boxtimes
%
\DeclareMathSymbol{\boxtimes}{\mathbin}{AMSa}{"02}

%
% Macros for moduli stacks/spaces
%
\def\QCohstack{\mathcal{QC}\!\mathit{oh}}
\def\Cohstack{\mathcal{C}\!\mathit{oh}}
\def\Spacesstack{\mathcal{S}\!\mathit{paces}}
\def\Quotfunctor{\mathrm{Quot}}
\def\Hilbfunctor{\mathrm{Hilb}}
\def\Curvesstack{\mathcal{C}\!\mathit{urves}}
\def\Polarizedstack{\mathcal{P}\!\mathit{olarized}}
\def\Complexesstack{\mathcal{C}\!\mathit{omplexes}}
% \Pic is the operator that assigns to X its picard group, usage \Pic(X)
% \Picardstack_{X/B} denotes the Picard stack of X over B
% \Picardfunctor_{X/B} denotes the Picard functor of X over B
\def\Pic{\mathop{\mathrm{Pic}}\nolimits}
\def\Picardstack{\mathcal{P}\!\mathit{ic}}
\def\Picardfunctor{\mathrm{Pic}}
\def\Deformationcategory{\mathcal{D}\!\mathit{ef}}
\newenvironment{reference}{}{}
\newenvironment{slogan}{}{}
\newenvironment{history}{}{}
\usepackage{cleveref}
\crefname{originallemma}{original lemma}{original lemmas}
\Crefname{originallemma}{Original lemma}{Original lemmas}
\makeatletter
\newlabel{cohomology-lemma-topology-I-adic-f}{{0EHA}{source}{}{sourcepointer-1}{}}
\newlabel{cohomology-lemma-topology-I-adic-f@cref}{{[originallemma][0][]0EHA}{[1][0][]source}{}{}{}}
\makeatother
\crefname{originallemma}{original lemma}{original lemmas}
\Crefname{originallemma}{Original lemma}{Original lemmas}
\makeatletter
\newlabel{local-cohomology-lemma-cd-is-one}{{0DXD}{source}{}{sourcepointer-2}{}}
\newlabel{local-cohomology-lemma-cd-is-one@cref}{{[originallemma][0][]0DXD}{[1][0][]source}{}{}{}}
\makeatother
\crefname{originallemma}{original lemma}{original lemmas}
\Crefname{originallemma}{Original lemma}{Original lemmas}
\makeatletter
\newlabel{coherent-lemma-homs-over-open}{{01YB}{source}{}{sourcepointer-3}{}}
\newlabel{coherent-lemma-homs-over-open@cref}{{[originallemma][0][]01YB}{[1][0][]source}{}{}{}}
\makeatother
\crefname{originalsection}{original section}{original sections}
\Crefname{originalsection}{Original section}{Original sections}
\makeatletter
\newlabel{coherent-section-coherent-formal}{{0EHN}{source}{}{sourcepointer-4}{}}
\newlabel{coherent-section-coherent-formal@cref}{{[originalsection][0][]0EHN}{[1][0][]source}{}{}{}}
\makeatother
\crefname{originallemma}{original lemma}{original lemmas}
\Crefname{originallemma}{Original lemma}{Original lemmas}
\makeatletter
\newlabel{coherent-lemma-inverse-systems-affine}{{087W}{source}{}{sourcepointer-5}{}}
\newlabel{coherent-lemma-inverse-systems-affine@cref}{{[originallemma][0][]087W}{[1][0][]source}{}{}{}}
\makeatother
\crefname{originallemma}{original lemma}{original lemmas}
\Crefname{originallemma}{Original lemma}{Original lemmas}
\makeatletter
\newlabel{lemma-properties-system}{{0EI8}{source}{}{sourcepointer-6}{}}
\newlabel{lemma-properties-system@cref}{{[originallemma][0][]0EI8}{[1][0][]source}{}{}{}}
\makeatother
\crefname{originallemma}{original lemma}{original lemmas}
\Crefname{originallemma}{Original lemma}{Original lemmas}
\makeatletter
\newlabel{algebra-lemma-finite-over-complete-ring}{{031D}{source}{}{sourcepointer-7}{}}
\newlabel{algebra-lemma-finite-over-complete-ring@cref}{{[originallemma][0][]031D}{[1][0][]source}{}{}{}}
\makeatother
\crefname{originallemma}{original lemma}{original lemmas}
\Crefname{originallemma}{Original lemma}{Original lemmas}
\makeatletter
\newlabel{local-cohomology-lemma-finiteness-pushforwards-and-H1-local}{{0BK0}{source}{}{sourcepointer-8}{}}
\newlabel{local-cohomology-lemma-finiteness-pushforwards-and-H1-local@cref}{{[originallemma][0][]0BK0}{[1][0][]source}{}{}{}}
\makeatother
\crefname{originallemma}{original lemma}{original lemmas}
\Crefname{originallemma}{Original lemma}{Original lemmas}
\makeatletter
\newlabel{local-cohomology-lemma-finiteness-for-finite-locally-free}{{0BLT}{source}{}{sourcepointer-9}{}}
\newlabel{local-cohomology-lemma-finiteness-for-finite-locally-free@cref}{{[originallemma][0][]0BLT}{[1][0][]source}{}{}{}}
\makeatother
\crefname{originallemma}{original lemma}{original lemmas}
\Crefname{originallemma}{Original lemma}{Original lemmas}
\makeatletter
\newlabel{cohomology-lemma-ML}{{0BLC}{source}{}{sourcepointer-10}{}}
\newlabel{cohomology-lemma-ML@cref}{{[originallemma][0][]0BLC}{[1][0][]source}{}{}{}}
\makeatother
\crefname{originaltheorem}{original theorem}{original theorems}
\Crefname{originaltheorem}{Original theorem}{Original theorems}
\makeatletter
\newlabel{local-cohomology-theorem-finiteness}{{0BJV}{source}{}{sourcepointer-11}{}}
\newlabel{local-cohomology-theorem-finiteness@cref}{{[originaltheorem][0][]0BJV}{[1][0][]source}{}{}{}}
\makeatother
\crefname{originallemma}{original lemma}{original lemmas}
\Crefname{originallemma}{Original lemma}{Original lemmas}
\makeatletter
\newlabel{lemma-algebraization-principal}{{0DXU}{source}{}{sourcepointer-12}{}}
\newlabel{lemma-algebraization-principal@cref}{{[originallemma][0][]0DXU}{[1][0][]source}{}{}{}}
\makeatother
\crefname{originalsection}{original section}{original sections}
\Crefname{originalsection}{Original section}{Original sections}
\makeatletter
\newlabel{remark-interesting-case}{{0DXV}{source}{}{sourcepointer-13}{}}
\newlabel{remark-interesting-case@cref}{{[originalsection][0][]0DXV}{[1][0][]source}{}{}{}}
\makeatother
\crefname{originallemma}{original lemma}{original lemmas}
\Crefname{originallemma}{Original lemma}{Original lemmas}
\makeatletter
\newlabel{local-cohomology-lemma-finiteness-Rjstar}{{0BJY}{source}{}{sourcepointer-14}{}}
\newlabel{local-cohomology-lemma-finiteness-Rjstar@cref}{{[originallemma][0][]0BJY}{[1][0][]source}{}{}{}}
\makeatother
\crefname{originallemma}{original lemma}{original lemmas}
\Crefname{originallemma}{Original lemma}{Original lemmas}
\makeatletter
\newlabel{cohomology-lemma-ML-general}{{0GYK}{source}{}{sourcepointer-15}{}}
\newlabel{cohomology-lemma-ML-general@cref}{{[originallemma][0][]0GYK}{[1][0][]source}{}{}{}}
\makeatother
\crefname{originalsection}{original section}{original sections}
\Crefname{originalsection}{Original section}{Original sections}
\makeatletter
\newlabel{section-completion}{{0EKN}{source}{}{sourcepointer-16}{}}
\newlabel{section-completion@cref}{{[originalsection][0][]0EKN}{[1][0][]source}{}{}{}}
\makeatother
\crefname{originallemma}{original lemma}{original lemmas}
\Crefname{originallemma}{Original lemma}{Original lemmas}
\makeatletter
\newlabel{dualizing-lemma-divide-by-torsion}{{0AW0}{source}{}{sourcepointer-17}{}}
\newlabel{dualizing-lemma-divide-by-torsion@cref}{{[originallemma][0][]0AW0}{[1][0][]source}{}{}{}}
\makeatother
\crefname{originalequation}{original equation}{original equations}
\Crefname{originalequation}{Original equation}{Original equations}
\makeatletter
\newlabel{coherent-equation-completion-functor}{{0880}{source}{}{sourcepointer-18}{}}
\newlabel{coherent-equation-completion-functor@cref}{{[originalequation][0][]0880}{[1][0][]source}{}{}{}}
\makeatother
\crefname{originallemma}{original lemma}{original lemmas}
\Crefname{originallemma}{Original lemma}{Original lemmas}
\makeatletter
\newlabel{lemma-fully-faithful-inequalities}{{0EJ7}{source}{}{sourcepointer-19}{}}
\newlabel{lemma-fully-faithful-inequalities@cref}{{[originallemma][0][]0EJ7}{[1][0][]source}{}{}{}}
\makeatother
\crefname{originallemma}{original lemma}{original lemmas}
\Crefname{originallemma}{Original lemma}{Original lemmas}
\makeatletter
\newlabel{properties-lemma-lift-finite-presentation}{{0G41}{source}{}{sourcepointer-20}{}}
\newlabel{properties-lemma-lift-finite-presentation@cref}{{[originallemma][0][]0G41}{[1][0][]source}{}{}{}}
\makeatother
\crefname{originallemma}{original lemma}{original lemmas}
\Crefname{originallemma}{Original lemma}{Original lemmas}
\makeatletter
\newlabel{divisors-lemma-extend-invertible-module}{{0BD9}{source}{}{sourcepointer-21}{}}
\newlabel{divisors-lemma-extend-invertible-module@cref}{{[originallemma][0][]0BD9}{[1][0][]source}{}{}{}}
\makeatother
\crefname{originallemma}{original lemma}{original lemmas}
\Crefname{originallemma}{Original lemma}{Original lemmas}
\makeatletter
\newlabel{algebra-lemma-completion-flat}{{00MB}{source}{}{sourcepointer-22}{}}
\newlabel{algebra-lemma-completion-flat@cref}{{[originallemma][0][]00MB}{[1][0][]source}{}{}{}}
\makeatother
\crefname{originallemma}{original lemma}{original lemmas}
\Crefname{originallemma}{Original lemma}{Original lemmas}
\makeatletter
\newlabel{algebra-lemma-completion-complete}{{031C}{source}{}{sourcepointer-23}{}}
\newlabel{algebra-lemma-completion-complete@cref}{{[originallemma][0][]031C}{[1][0][]source}{}{}{}}
\makeatother
\crefname{originallemma}{original lemma}{original lemmas}
\Crefname{originallemma}{Original lemma}{Original lemmas}
\makeatletter
\newlabel{coherent-lemma-inverse-systems-pullback-equivalence}{{0EHQ}{source}{}{sourcepointer-24}{}}
\newlabel{coherent-lemma-inverse-systems-pullback-equivalence@cref}{{[originallemma][0][]0EHQ}{[1][0][]source}{}{}{}}
\makeatother
\crefname{originallemma}{original lemma}{original lemmas}
\Crefname{originallemma}{Original lemma}{Original lemmas}
\makeatletter
\newlabel{lemma-recognize-formal-coherent-modules}{{0EIQ}{source}{}{sourcepointer-25}{}}
\newlabel{lemma-recognize-formal-coherent-modules@cref}{{[originallemma][0][]0EIQ}{[1][0][]source}{}{}{}}
\makeatother
\crefname{originallemma}{original lemma}{original lemmas}
\Crefname{originallemma}{Original lemma}{Original lemmas}
\makeatletter
\newlabel{properties-lemma-codimension-local-ring}{{02IZ}{source}{}{sourcepointer-26}{}}
\newlabel{properties-lemma-codimension-local-ring@cref}{{[originallemma][0][]02IZ}{[1][0][]source}{}{}{}}
\makeatother
\crefname{originallemma}{original lemma}{original lemmas}
\Crefname{originallemma}{Original lemma}{Original lemmas}
\makeatletter
\newlabel{morphisms-lemma-dimension-formula}{{02JU}{source}{}{sourcepointer-27}{}}
\newlabel{morphisms-lemma-dimension-formula@cref}{{[originallemma][0][]02JU}{[1][0][]source}{}{}{}}
\makeatother
\crefname{originallemma}{original lemma}{original lemmas}
\Crefname{originallemma}{Original lemma}{Original lemmas}
\makeatletter
\newlabel{properties-lemma-quasi-coherent-quasi-affine}{{0EHM}{source}{}{sourcepointer-28}{}}
\newlabel{properties-lemma-quasi-coherent-quasi-affine@cref}{{[originallemma][0][]0EHM}{[1][0][]source}{}{}{}}
\makeatother
\crefname{originallemma}{original lemma}{original lemmas}
\Crefname{originallemma}{Original lemma}{Original lemmas}
\makeatletter
\newlabel{algebra-lemma-limit-complete}{{09B8}{source}{}{sourcepointer-29}{}}
\newlabel{algebra-lemma-limit-complete@cref}{{[originallemma][0][]09B8}{[1][0][]source}{}{}{}}
\makeatother
\crefname{originallemma}{original lemma}{original lemmas}
\Crefname{originallemma}{Original lemma}{Original lemmas}
\makeatletter
\newlabel{algebra-lemma-hathat-finitely-generated}{{05GG}{source}{}{sourcepointer-30}{}}
\newlabel{algebra-lemma-hathat-finitely-generated@cref}{{[originallemma][0][]05GG}{[1][0][]source}{}{}{}}
\makeatother
\title{Canonical algebraization of coherent formal systems under cohomological-dimension-one depth bounds}
\author{The Stacks Project Authors}
\date{stacks-project:algebraization.tex:local-snapshot-20261011; GNU FDL 1.2 or later}
\begin{document}
\maketitle
\noindent\textit{Editorial note.} The exactcd(A,I)=1 assumption and the depth-plus-distance condition on a pro-isomorphic coherent system remain in the statement. The principal-ideal algebraization variant is not counted separately.\par\medskip
\noindent\textit{Allowed standard prerequisites (editorial).}\begin{itemize}
\item Explicit prerequisite boundary. Dualizing complexes, coherent formal modules, local-cohomology finiteness from depth, adic completion, finite modules over complete rings and Mittag–Leffler limits. The source-local topology and canonical-extension criterion are retained.
\end{itemize}\medskip
\ifdefined\hypertarget\hypertarget{sourcepointer-1}{}\fi
\noindent\textit{Source reference.} Original lemma 0EHA: focused-seeds/AGP/sources/stacks-snapshot-20261011/tags.txt, source lines 18125--18125. Stable Stacks tag for an expressly allowed foundational or notational input; the counted construction remains the two concrete core obligations. The retained local core is not replaced by this pointer. See standard\_prerequisites.\par\medskip
\ifdefined\hypertarget\hypertarget{sourcepointer-2}{}\fi
\noindent\textit{Source reference.} Original lemma 0DXD: focused-seeds/AGP/sources/stacks-snapshot-20261011/tags.txt, source lines 17428--17428. Stable Stacks tag for an expressly allowed foundational or notational input; the counted construction remains the two concrete core obligations. The retained local core is not replaced by this pointer. See standard\_prerequisites.\par\medskip
\ifdefined\hypertarget\hypertarget{sourcepointer-3}{}\fi
\noindent\textit{Source reference.} Original lemma 01YB: focused-seeds/AGP/sources/stacks-snapshot-20261011/tags.txt, source lines 2550--2550. Stable Stacks tag for an expressly allowed foundational or notational input; the counted construction remains the two concrete core obligations. The retained local core is not replaced by this pointer. See standard\_prerequisites.\par\medskip
\ifdefined\hypertarget\hypertarget{sourcepointer-4}{}\fi
\noindent\textit{Source reference.} Original section 0EHN: focused-seeds/AGP/sources/stacks-snapshot-20261011/tags.txt, source lines 18138--18138. Stable Stacks tag for an expressly allowed foundational or notational input; the counted construction remains the two concrete core obligations. The retained local core is not replaced by this pointer. See standard\_prerequisites.\par\medskip
\ifdefined\hypertarget\hypertarget{sourcepointer-5}{}\fi
\noindent\textit{Source reference.} Original lemma 087W: focused-seeds/AGP/sources/stacks-snapshot-20261011/tags.txt, source lines 10441--10441. Stable Stacks tag for an expressly allowed foundational or notational input; the counted construction remains the two concrete core obligations. The retained local core is not replaced by this pointer. See standard\_prerequisites.\par\medskip
\ifdefined\hypertarget\hypertarget{sourcepointer-6}{}\fi
\noindent\textit{Source reference.} Original lemma 0EI8: focused-seeds/AGP/sources/stacks-snapshot-20261011/tags.txt, source lines 18158--18158. Stable Stacks tag for an expressly allowed foundational or notational input; the counted construction remains the two concrete core obligations. The retained local core is not replaced by this pointer. See standard\_prerequisites.\par\medskip
\ifdefined\hypertarget\hypertarget{sourcepointer-7}{}\fi
\noindent\textit{Source reference.} Original lemma 031D: focused-seeds/AGP/sources/stacks-snapshot-20261011/tags.txt, source lines 3980--3980. Stable Stacks tag for an expressly allowed foundational or notational input; the counted construction remains the two concrete core obligations. The retained local core is not replaced by this pointer. See standard\_prerequisites.\par\medskip
\ifdefined\hypertarget\hypertarget{sourcepointer-8}{}\fi
\noindent\textit{Source reference.} Original lemma 0BK0: focused-seeds/AGP/sources/stacks-snapshot-20261011/tags.txt, source lines 14543--14543. Stable Stacks tag for an expressly allowed foundational or notational input; the counted construction remains the two concrete core obligations. The retained local core is not replaced by this pointer. See standard\_prerequisites.\par\medskip
\ifdefined\hypertarget\hypertarget{sourcepointer-9}{}\fi
\noindent\textit{Source reference.} Original lemma 0BLT: focused-seeds/AGP/sources/stacks-snapshot-20261011/tags.txt, source lines 14606--14606. Stable Stacks tag for an expressly allowed foundational or notational input; the counted construction remains the two concrete core obligations. The retained local core is not replaced by this pointer. See standard\_prerequisites.\par\medskip
\ifdefined\hypertarget\hypertarget{sourcepointer-10}{}\fi
\noindent\textit{Source reference.} Original lemma 0BLC: focused-seeds/AGP/sources/stacks-snapshot-20261011/tags.txt, source lines 14590--14590. Stable Stacks tag for an expressly allowed foundational or notational input; the counted construction remains the two concrete core obligations. The retained local core is not replaced by this pointer. See standard\_prerequisites.\par\medskip
\ifdefined\hypertarget\hypertarget{sourcepointer-11}{}\fi
\noindent\textit{Source reference.} Original theorem 0BJV: focused-seeds/AGP/sources/stacks-snapshot-20261011/tags.txt, source lines 14538--14538. Stable Stacks tag for an expressly allowed foundational or notational input; the counted construction remains the two concrete core obligations. The retained local core is not replaced by this pointer. See standard\_prerequisites.\par\medskip
\ifdefined\hypertarget\hypertarget{sourcepointer-12}{}\fi
\noindent\textit{Source reference.} Original lemma 0DXU: focused-seeds/AGP/sources/stacks-snapshot-20261011/tags.txt, source lines 17444--17444. Stable Stacks tag for an expressly allowed foundational or notational input; the counted construction remains the two concrete core obligations. The retained local core is not replaced by this pointer. See standard\_prerequisites.\par\medskip
\ifdefined\hypertarget\hypertarget{sourcepointer-13}{}\fi
\noindent\textit{Source reference.} Original section 0DXV: focused-seeds/AGP/sources/stacks-snapshot-20261011/tags.txt, source lines 17445--17445. Stable Stacks tag for an expressly allowed foundational or notational input; the counted construction remains the two concrete core obligations. The retained local core is not replaced by this pointer. See standard\_prerequisites.\par\medskip
\ifdefined\hypertarget\hypertarget{sourcepointer-14}{}\fi
\noindent\textit{Source reference.} Original lemma 0BJY: focused-seeds/AGP/sources/stacks-snapshot-20261011/tags.txt, source lines 14541--14541. Stable Stacks tag for an expressly allowed foundational or notational input; the counted construction remains the two concrete core obligations. The retained local core is not replaced by this pointer. See standard\_prerequisites.\par\medskip
\ifdefined\hypertarget\hypertarget{sourcepointer-15}{}\fi
\noindent\textit{Source reference.} Original lemma 0GYK: focused-seeds/AGP/sources/stacks-snapshot-20261011/tags.txt, source lines 21145--21145. Stable Stacks tag for an expressly allowed foundational or notational input; the counted construction remains the two concrete core obligations. The retained local core is not replaced by this pointer. See standard\_prerequisites.\par\medskip
\ifdefined\hypertarget\hypertarget{sourcepointer-16}{}\fi
\noindent\textit{Source reference.} Original section 0EKN: focused-seeds/AGP/sources/stacks-snapshot-20261011/tags.txt, source lines 18243--18243. Stable Stacks tag for an expressly allowed foundational or notational input; the counted construction remains the two concrete core obligations. The retained local core is not replaced by this pointer. See standard\_prerequisites.\par\medskip
\ifdefined\hypertarget\hypertarget{sourcepointer-17}{}\fi
\noindent\textit{Source reference.} Original lemma 0AW0: focused-seeds/AGP/sources/stacks-snapshot-20261011/tags.txt, source lines 13703--13703. Stable Stacks tag for an expressly allowed foundational or notational input; the counted construction remains the two concrete core obligations. The retained local core is not replaced by this pointer. See standard\_prerequisites.\par\medskip
\ifdefined\hypertarget\hypertarget{sourcepointer-18}{}\fi
\noindent\textit{Source reference.} Original equation 0880: focused-seeds/AGP/sources/stacks-snapshot-20261011/tags.txt, source lines 10445--10445. Stable Stacks tag for an expressly allowed foundational or notational input; the counted construction remains the two concrete core obligations. The retained local core is not replaced by this pointer. See standard\_prerequisites.\par\medskip
\ifdefined\hypertarget\hypertarget{sourcepointer-19}{}\fi
\noindent\textit{Source reference.} Original lemma 0EJ7: focused-seeds/AGP/sources/stacks-snapshot-20261011/tags.txt, source lines 18192--18192. Stable Stacks tag for an expressly allowed foundational or notational input; the counted construction remains the two concrete core obligations. The retained local core is not replaced by this pointer. See standard\_prerequisites.\par\medskip
\ifdefined\hypertarget\hypertarget{sourcepointer-20}{}\fi
\noindent\textit{Source reference.} Original lemma 0G41: focused-seeds/AGP/sources/stacks-snapshot-20261011/tags.txt, source lines 20111--20111. Stable Stacks tag for an expressly allowed foundational or notational input; the counted construction remains the two concrete core obligations. The retained local core is not replaced by this pointer. See standard\_prerequisites.\par\medskip
\ifdefined\hypertarget\hypertarget{sourcepointer-21}{}\fi
\noindent\textit{Source reference.} Original lemma 0BD9: focused-seeds/AGP/sources/stacks-snapshot-20261011/tags.txt, source lines 14307--14307. Stable Stacks tag for an expressly allowed foundational or notational input; the counted construction remains the two concrete core obligations. The retained local core is not replaced by this pointer. See standard\_prerequisites.\par\medskip
\ifdefined\hypertarget\hypertarget{sourcepointer-22}{}\fi
\noindent\textit{Source reference.} Original lemma 00MB: focused-seeds/AGP/sources/stacks-snapshot-20261011/tags.txt, source lines 817--817. Stable Stacks tag for an expressly allowed foundational or notational input; the counted construction remains the two concrete core obligations. The retained local core is not replaced by this pointer. See standard\_prerequisites.\par\medskip
\ifdefined\hypertarget\hypertarget{sourcepointer-23}{}\fi
\noindent\textit{Source reference.} Original lemma 031C: focused-seeds/AGP/sources/stacks-snapshot-20261011/tags.txt, source lines 3979--3979. Stable Stacks tag for an expressly allowed foundational or notational input; the counted construction remains the two concrete core obligations. The retained local core is not replaced by this pointer. See standard\_prerequisites.\par\medskip
\ifdefined\hypertarget\hypertarget{sourcepointer-24}{}\fi
\noindent\textit{Source reference.} Original lemma 0EHQ: focused-seeds/AGP/sources/stacks-snapshot-20261011/tags.txt, source lines 18140--18140. Stable Stacks tag for an expressly allowed foundational or notational input; the counted construction remains the two concrete core obligations. The retained local core is not replaced by this pointer. See standard\_prerequisites.\par\medskip
\ifdefined\hypertarget\hypertarget{sourcepointer-25}{}\fi
\noindent\textit{Source reference.} Original lemma 0EIQ: focused-seeds/AGP/sources/stacks-snapshot-20261011/tags.txt, source lines 18175--18175. Stable Stacks tag for an expressly allowed foundational or notational input; the counted construction remains the two concrete core obligations. The retained local core is not replaced by this pointer. See standard\_prerequisites.\par\medskip
\ifdefined\hypertarget\hypertarget{sourcepointer-26}{}\fi
\noindent\textit{Source reference.} Original lemma 02IZ: focused-seeds/AGP/sources/stacks-snapshot-20261011/tags.txt, source lines 3318--3318. Stable Stacks tag for an expressly allowed foundational or notational input; the counted construction remains the two concrete core obligations. The retained local core is not replaced by this pointer. See standard\_prerequisites.\par\medskip
\ifdefined\hypertarget\hypertarget{sourcepointer-27}{}\fi
\noindent\textit{Source reference.} Original lemma 02JU: focused-seeds/AGP/sources/stacks-snapshot-20261011/tags.txt, source lines 3349--3349. Stable Stacks tag for an expressly allowed foundational or notational input; the counted construction remains the two concrete core obligations. The retained local core is not replaced by this pointer. See standard\_prerequisites.\par\medskip
\ifdefined\hypertarget\hypertarget{sourcepointer-28}{}\fi
\noindent\textit{Source reference.} Original lemma 0EHM: focused-seeds/AGP/sources/stacks-snapshot-20261011/tags.txt, source lines 18137--18137. Classical finite-module completion or quasi-affine sheaf fact in the completed context proof; not the canonical-extension construction.\par\medskip
\ifdefined\hypertarget\hypertarget{sourcepointer-29}{}\fi
\noindent\textit{Source reference.} Original lemma 09B8: focused-seeds/AGP/sources/stacks-snapshot-20261011/tags.txt, source lines 11786--11786. Classical finite-module completion or quasi-affine sheaf fact in the completed context proof; not the canonical-extension construction.\par\medskip
\ifdefined\hypertarget\hypertarget{sourcepointer-30}{}\fi
\noindent\textit{Source reference.} Original lemma 05GG: focused-seeds/AGP/sources/stacks-snapshot-20261011/tags.txt, source lines 7060--7060. Classical finite-module completion or quasi-affine sheaf fact in the completed context proof; not the canonical-extension construction.\par\medskip
\begin{lemma}
\label{lemma-cd-1-canonical}
In Situation \ref{situation-algebraize} let $(\mathcal{F}_n)$
be an object of $\textit{Coh}(U, I\mathcal{O}_U)$. Assume
\begin{enumerate}
\item $A$ has a dualizing complex and $\text{cd}(A, I) = 1$,
\item $(\mathcal{F}_n)$ is pro-isomorphic to an inverse system
$(\mathcal{F}_n'')$ of coherent $\mathcal{O}_U$-modules such that
$\text{depth}(\mathcal{F}''_{n, y}) + \delta^Y_Z(y) \geq 3$
for all $y \in U \cap Y$.
\end{enumerate}
Then $(\mathcal{F}_n)$ extends canonically to $X$, see
Definition \ref{definition-canonically-algebraizable}.
\end{lemma}

\begin{situation}
\label{situation-algebraize}
Here $A$ is a Noetherian ring and $I \subset \mathfrak a \subset A$ are ideals.
We set $X = \Spec(A)$, $Y = V(I) = \Spec(A/I)$, and
$Z = V(\mathfrak a) = \Spec(A/\mathfrak a)$. Furthermore $U = X \setminus Z$.
\end{situation}

\noindent
In this section we try to find conditions that guarantee an object
of $\textit{Coh}(U, I\mathcal{O}_U)$ is in the image of the completion functor
$\textit{Coh}(\mathcal{O}_U) \to \textit{Coh}(U, I\mathcal{O}_U)$.
See Cohomology of Schemes, Section \ref{coherent-section-coherent-formal} and
Section \ref{section-completion}.

\begin{lemma}
\label{lemma-system-of-modules}
In Situation \ref{situation-algebraize}.
Consider an inverse system $(M_n)$ of $A$-modules such
that
\begin{enumerate}
\item $M_n$ is a finite $A$-module,
\item $M_n$ is annihilated by $I^n$,
\item the kernel and cokernel of $M_{n + 1}/I^nM_{n + 1} \to M_n$
are $\mathfrak a$-power torsion.
\end{enumerate}
Then $(\widetilde{M}_n|_U)$ is in $\textit{Coh}(U, I\mathcal{O}_U)$.
Conversely, every object of $\textit{Coh}(U, I\mathcal{O}_U)$
arises in this manner.
\end{lemma}

\begin{proof}
We omit the verification that $(\widetilde{M}_n|_U)$ is in
$\textit{Coh}(U, I\mathcal{O}_U)$. Let $(\mathcal{F}_n)$
be an object of $\textit{Coh}(U, I\mathcal{O}_U)$.
By Local Cohomology, Lemma
\ref{local-cohomology-lemma-finiteness-pushforwards-and-H1-local}
we see that $\mathcal{F}_n = \widetilde{M_n}$ for some finite
$A/I^n$-module $M_n$. After dividing $M_n$ by $H^0_\mathfrak a(M_n)$
we may assume $M_n \subset H^0(U, \mathcal{F}_n)$, see
Dualizing Complexes, Lemma \ref{dualizing-lemma-divide-by-torsion}
and the already referenced lemma.
After replacing inductively $M_{n + 1}$ by the inverse image
of $M_n$ under the map $M_{n + 1} \to H^0(U, \mathcal{F}_{n + 1})
\to H^0(U, \mathcal{F}_n)$, we may assume $M_{n + 1}$ maps into
$M_n$. This gives an inverse system $(M_n)$ satisfying (1) and (2)
such that $\mathcal{F}_n = \widetilde{M_n}$. To see that (3)
holds, use that $M_{n + 1}/I^nM_{n + 1} \to M_n$ is a map
of finite $A$-modules which induces an isomorphism after
applying $\widetilde{\ }$ and restriction to $U$
(here we use the first referenced lemma one more time).
\end{proof}

\noindent
In Situation \ref{situation-algebraize} we can study the completion functor
Cohomology of Schemes, Equation (\ref{coherent-equation-completion-functor})
\begin{equation}
\label{equation-completion}
\textit{Coh}(\mathcal{O}_U)
\longrightarrow
\textit{Coh}(U, I\mathcal{O}_U),\quad
\mathcal{F} \longmapsto \mathcal{F}^\wedge
\end{equation}
If $A$ is $I$-adically complete, then this functor is fully faithful
on suitable subcategories by our earlier work on algebraization of
formal sections, see Section \ref{section-completion} and
Lemma \ref{lemma-fully-faithful-inequalities} for some sample results.
Next, let $(\mathcal{F}_n)$ be an object of $\textit{Coh}(U, I\mathcal{O}_U)$.
Still assuming $A$ is $I$-adically complete, we can ask:
When is $(\mathcal{F}_n)$ in the essential image of the completion
functor displayed above?

\begin{lemma}
\label{lemma-essential-image-completion}
In Situation \ref{situation-algebraize} let $(\mathcal{F}_n)$
be an object of $\textit{Coh}(U, I\mathcal{O}_U)$. Consider the
following conditions:
\begin{enumerate}
\item $(\mathcal{F}_n)$ is in the essential image
of the functor (\ref{equation-completion}),
\item $(\mathcal{F}_n)$ is the completion of a
coherent $\mathcal{O}_U$-module,
\item $(\mathcal{F}_n)$ is the completion of a coherent
$\mathcal{O}_V$-module for $U \cap Y \subset V \subset U$ open,
\item $(\mathcal{F}_n)$ is the completion of
the restriction to $U$ of a coherent $\mathcal{O}_X$-module,
\item $(\mathcal{F}_n)$ is the restriction to $U$ of
the completion of a coherent $\mathcal{O}_X$-module,
\item there exists an object $(\mathcal{G}_n)$ of
$\textit{Coh}(X, I\mathcal{O}_X)$ whose restriction
to $U$ is $(\mathcal{F}_n)$.
\end{enumerate}
Then conditions (1), (2), (3), (4), and (5) are equivalent and imply (6).
If $A$ is $I$-adically complete then condition (6) implies the others.
\end{lemma}

\begin{proof}
Parts (1) and (2) are equivalent, because the completion of a coherent
$\mathcal{O}_U$-module $\mathcal{F}$ is by definition the image of
$\mathcal{F}$ under the functor (\ref{equation-completion}).
If $V \subset U$ is an open subscheme containing $U \cap Y$, then we have
$$
\textit{Coh}(V, I\mathcal{O}_V) =
\textit{Coh}(U, I\mathcal{O}_U)
$$
since the category of coherent $\mathcal{O}_V$-modules supported on
$V \cap Y$ is the same as the category of coherent $\mathcal{O}_U$-modules
supported on $U \cap Y$. Thus the completion of a coherent
$\mathcal{O}_V$-module is an object of $\textit{Coh}(U, I\mathcal{O}_U)$.
Having said this the equivalence of (2), (3), (4), and (5)
holds because the functors
$\textit{Coh}(\mathcal{O}_X) \to \textit{Coh}(\mathcal{O}_U) \to
\textit{Coh}(\mathcal{O}_V)$ are essentially surjective.
See Properties, Lemma \ref{properties-lemma-lift-finite-presentation}.

\medskip\noindent
It is always the case that (5) implies (6). Assume $A$ is $I$-adically complete.
Then any object of $\textit{Coh}(X, I\mathcal{O}_X)$ corresponds to a finite
$A$-module by Cohomology of Schemes, Lemma
\ref{coherent-lemma-inverse-systems-affine}.
Thus we see that (6) implies (5) in this case.
\end{proof}

\begin{example}
\label{example-not-algebraizable}
Let $k$ be a field. Let $A = k[x, y][[t]]$ with $I = (t)$ and
$\mathfrak a = (x, y, t)$. Let us use notation as in
Situation \ref{situation-algebraize}. Observe that
$U \cap Y = (D(x) \cap Y) \cup (D(y) \cap Y)$ is an affine
open covering. For $n \geq 1$ consider the invertible
module $\mathcal{L}_n$ of $\mathcal{O}_U/t^n\mathcal{O}_U$
given by glueing $A_x/t^nA_x$ and $A_y/t^nA_y$ via the invertible
element of $A_{xy}/t^nA_{xy}$ which is the image of any power series
of the form
$$
u = 1 + \frac{t}{xy} + \sum_{n \geq 2} a_n \frac{t^n}{(xy)^{\varphi(n)}}
$$
with $a_n \in k[x, y]$ and $\varphi(n) \in \mathbf{N}$.
Then $(\mathcal{L}_n)$ is an invertible object of
$\textit{Coh}(U, I\mathcal{O}_U)$ which is not the
completion of a coherent $\mathcal{O}_U$-module $\mathcal{L}$.
We only sketch the argument and we omit most of the details.
Let $y \in U \cap Y$. Then the completion of the stalk
$\mathcal{L}_y$ would be an invertible module hence $\mathcal{L}_y$
is invertible. Thus there would exist an open $V \subset U$
containing $U \cap Y$ such that $\mathcal{L}|_V$ is invertible.
By Divisors, Lemma \ref{divisors-lemma-extend-invertible-module}
we find an invertible $A$-module $M$ with
$\widetilde{M}|_V \cong \mathcal{L}|_V$. However the ring $A$
is a UFD hence we see $M \cong A$ which would imply
$\mathcal{L}_n \cong \mathcal{O}_U/I^n\mathcal{O}_U$.
Since $\mathcal{L}_2 \not \cong \mathcal{O}_U/I^2\mathcal{O}_U$
by construction we get a contradiction as desired.

\medskip\noindent
Note that if we take $a_n = 0$ for $n \geq 2$, then we see
that $\lim H^0(U, \mathcal{L}_n)$ is nonzero: in this case
we the function $x$ on $D(x)$ and the function $x + t/y$ on $D(y)$
glue. On the other hand, if we take $a_n = 1$ and $\varphi(n) = 2^n$
or even $\varphi(n) = n^2$ then the reader can show that
$\lim H^0(U, \mathcal{L}_n)$ is zero; this gives another proof
that $(\mathcal{L}_n)$ is not algebraizable in this case.
\end{example}

\noindent
If in Situation \ref{situation-algebraize} the ring $A$ is not
$I$-adically complete, then Lemma \ref{lemma-essential-image-completion}
suggests the correct thing is to ask whether $(\mathcal{F}_n)$
is in the essential image of the restriction functor
$$
\textit{Coh}(X, I\mathcal{O}_X)
\longrightarrow
\textit{Coh}(U, I\mathcal{O}_U)
$$
However, we can no longer say that this means $(\mathcal{F}_n)$
is algebraizable. Thus we introduce the following terminology.

\begin{definition}
\label{definition-algebraizable}
In Situation \ref{situation-algebraize} let $(\mathcal{F}_n)$ be an
object of $\textit{Coh}(U, I\mathcal{O}_U)$. We say
{\it $(\mathcal{F}_n)$ extends to $X$} if there exists an object
$(\mathcal{G}_n)$ of $\textit{Coh}(X, I\mathcal{O}_X)$ whose restriction
to $U$ is isomorphic to $(\mathcal{F}_n)$.
\end{definition}

\noindent
This notion is equivalent to being algebraizable over the completion.

\begin{lemma}
\label{lemma-algebraizable}
In Situation \ref{situation-algebraize} let $(\mathcal{F}_n)$ be an
object of $\textit{Coh}(U, I\mathcal{O}_U)$. Let $A', I', \mathfrak a'$
be the $I$-adic completions of $A, I, \mathfrak a$. Set $X' = \Spec(A')$
and $U' = X' \setminus V(\mathfrak a')$. The following are equivalent
\begin{enumerate}
\item $(\mathcal{F}_n)$ extends to $X$, and
\item the pullback of $(\mathcal{F}_n)$ to $U'$ is the completion
of a coherent $\mathcal{O}_{U'}$-module.
\end{enumerate}
\end{lemma}

\begin{proof}
Recall that $A \to A'$ is a flat ring map which induces an isomorphism
$A/I \to A'/I'$. See
Algebra, Lemmas \ref{algebra-lemma-completion-flat} and
\ref{algebra-lemma-completion-complete}.
Thus $X' \to X$ is a
flat morphism inducing an isomorphism $Y' \to Y$. Thus $U' \to U$
is a flat morphism which induces an isomorphism $U' \cap Y' \to U \cap Y$.
This implies that in the commutative diagram
$$
\xymatrix{
\textit{Coh}(X', I\mathcal{O}_{X'}) \ar[r] &
\textit{Coh}(U', I\mathcal{O}_{U'}) \\
\textit{Coh}(X, I\mathcal{O}_X) \ar[u] \ar[r] &
\textit{Coh}(U, I\mathcal{O}_U) \ar[u]
}
$$
the vertical functors are equivalences. See
Cohomology of Schemes, Lemma
\ref{coherent-lemma-inverse-systems-pullback-equivalence}.
The lemma follows formally from this and the results of
Lemma \ref{lemma-essential-image-completion}.
\end{proof}

\noindent
In Situation \ref{situation-algebraize} let $(\mathcal{F}_n)$ be an
object of $\textit{Coh}(U, I\mathcal{O}_U)$. To figure out if
$(\mathcal{F}_n)$ extends to $X$ it makes sense to look at the $A$-module
\begin{equation}
\label{equation-guess}
M = \lim H^0(U, \mathcal{F}_n)
\end{equation}
Observe that $M$ has a limit topology which is (a priori) coarser than
the $I$-adic topology since $M \to H^0(U, \mathcal{F}_n)$ annihilates $I^nM$.
There are canonical maps
$$
\widetilde{M}|_U \to \widetilde{M/I^nM}|_U \to
\widetilde{H^0(U, \mathcal{F}_n)}|_U \to
\mathcal{F}_n
$$
One could hope that $\widetilde{M}$ restricts to a coherent module
on $U$ and that $(\mathcal{F}_n)$ is the completion of this module.
This is naive because this has almost no chance of being true
if $A$ is not complete. But even if $A$ is $I$-adically complete
this notion is very difficult to work with.
A less naive approach is to consider the following requirement.

\begin{definition}
\label{definition-canonically-algebraizable}
In Situation \ref{situation-algebraize} let $(\mathcal{F}_n)$ be an
object of $\textit{Coh}(U, I\mathcal{O}_U)$. We say
{\it $(\mathcal{F}_n)$ canonically extends to $X$} if the the
inverse system
$$
\{\widetilde{H^0(U, \mathcal{F}_n)}\}_{n \geq 1}
$$
in $\QCoh(\mathcal{O}_X)$ is pro-isomorphic to an object
$(\mathcal{G}_n)$ of $\textit{Coh}(X, I\mathcal{O}_X)$.
\end{definition}

\noindent
We will see in Lemma \ref{lemma-canonically-algebraizable}
that the condition in Definition \ref{definition-canonically-algebraizable}
is stronger than the condition of Definition \ref{definition-algebraizable}.

\begin{lemma}
\label{lemma-canonically-algebraizable}
In Situation \ref{situation-algebraize} let $(\mathcal{F}_n)$ be an
object of $\textit{Coh}(U, I\mathcal{O}_U)$. If $(\mathcal{F}_n)$
canonically extends to $X$, then
\begin{enumerate}
\item $(\widetilde{H^0(U, \mathcal{F}_n)})$ is pro-isomorphic to an
object $(\mathcal{G}_n)$ of $\textit{Coh}(X, I \mathcal{O}_X)$
unique up to unique isomorphism,
\item the restriction of $(\mathcal{G}_n)$ to $U$ is isomorphic
to $(\mathcal{F}_n)$, i.e., $(\mathcal{F}_n)$ extends to $X$,
\item the inverse system $\{H^0(U, \mathcal{F}_n)\}$
satisfies the Mittag-Leffler condition, and
\item the module $M$ in (\ref{equation-guess}) is finite over the
$I$-adic completion of $A$ and the limit topology on
$M$ is the $I$-adic topology.
\end{enumerate}
\end{lemma}

\begin{proof}
The existence of $(\mathcal{G}_n)$ in (1) follows from
Definition \ref{definition-canonically-algebraizable}.
The uniqueness of $(\mathcal{G}_n)$ in (1) follows from
Lemma \ref{lemma-recognize-formal-coherent-modules}.
Write $\mathcal{G}_n = \widetilde{M_n}$.
Then $\{M_n\}$ is an inverse system of finite $A$-modules
with $M_n = M_{n + 1}/I^n M_{n + 1}$.
By Definition \ref{definition-canonically-algebraizable}
the inverse system $\{H^0(U, \mathcal{F}_n)\}$
is pro-isomorphic to $\{M_n\}$.
Hence we see that the inverse system $\{H^0(U, \mathcal{F}_n)\}$
satisfies the Mittag-Leffler condition and that
$M = \lim M_n$ (as topological modules).
Thus the properties of $M$ in (4) follow from
Algebra, Lemmas \ref{algebra-lemma-limit-complete},
\ref{algebra-lemma-finite-over-complete-ring}, and
\ref{algebra-lemma-hathat-finitely-generated}.
Since $U$ is quasi-affine the canonical maps
$$
\widetilde{H^0(U, \mathcal{F}_n)}|_U \to \mathcal{F}_n
$$
are isomorphisms (Properties, Lemma
\ref{properties-lemma-quasi-coherent-quasi-affine}).
We conclude that $(\mathcal{G}_n|_U)$ and $(\mathcal{F}_n)$ are
pro-isomorphic and hence isomorphic by
Lemma \ref{lemma-recognize-formal-coherent-modules}.
\end{proof}

\section{A distance function}
\label{section-distance}

\noindent
Let $Y$ be a Noetherian scheme and let $Z \subset Y$ be a closed subset.
We define a function
\begin{equation}
\label{equation-delta-Z}
\delta^Y_Z = \delta_Z : Y \longrightarrow \mathbf{Z}_{\geq 0} \cup \{\infty\}
\end{equation}
which measures the ``distance'' of a point of $Y$ from $Z$.
For an informal discussion, please see Remark \ref{remark-discussion}.
Let $y \in Y$. We set $\delta_Z(y) = \infty$ if $y$ is contained
in a connected component of $Y$ which does not meet $Z$.
If $y$ is contained in a connected component of $Y$ which meets $Z$,
then we can find $k \geq 0$ and a system
$$
V_0 \subset W_0 \supset V_1 \subset W_1 \supset \ldots \supset
V_k \subset W_k
$$
of integral closed subschemes of $Y$ such that $V_0 \subset Z$
and $y \in W_k$ is the generic point. Set
$c_i = \text{codim}(V_i, W_i)$ for $i = 0, \ldots, k$
and $b_i = \text{codim}(V_{i + 1}, W_i)$ for $i = 0, \ldots, k - 1$.
For such a system we set
$$
\delta(V_0, W_0, V_1, \ldots, W_k) =
k +
\max_{i = 0, 1, \ldots, k}
(c_i + c_{i + 1} + \ldots + c_k - b_i - b_{i + 1} - \ldots - b_{k - 1})
$$
This is $\geq k$ as we can take $i = k$ and we have $c_k \geq 0$.
Finally, we set
$$
\delta_Z(y) = \min \delta(V_0, W_0, V_1, \ldots, W_k)
$$
where the minimum is over all systems of integral closed subschemes of $Y$
as above.

\begin{lemma}
\label{lemma-discussion}
Let $Y$ be a Noetherian scheme and let $Z \subset Y$ be a closed subset.
\begin{enumerate}
\item For $y \in Y$ we have $\delta_Z(y) = 0 \Leftrightarrow y \in Z$.
\item The subsets $\{y \in Y \mid \delta_Z(y) \leq k\}$ are
stable under specialization.
\item For $y \in Y$ and $z \in \overline{\{y\}} \cap Z$ we have
$\dim(\mathcal{O}_{\overline{\{y\}}, z}) \geq \delta_Z(y)$.
\item If $\delta$ is a dimension function on $Y$, then
$\delta_Z(y) \geq \delta(y) - \delta_{max}$ where $\delta_{max}$
is the maximum value of $\delta$ on $Z$.
\item If $Y = \Spec(A)$ is the spectrum of a catenary Noetherian local ring
with maximal ideal $\mathfrak m$ and $Z = \{\mathfrak m\}$, then
$\delta_Z(y) = \dim(\overline{\{y\}})$.
\item If $Y' \subset Y$ is an open subscheme, then
$\delta^{Y'}_{Y' \cap Z}(y') \geq \delta^Y_Z(y')$ for $y' \in Y'$.
\end{enumerate}
Assume $Y$ is catenary. Then
\begin{enumerate}
\setcounter{enumi}{6}
\item Let $y' \leadsto y$ be an immediate specialization of points of $Y$.
If $Y$ is catenary, then $\delta_Z(y') \leq \delta_Z(y) + 1$.
\item Given a pattern of specializations
$$
\xymatrix{
& y'_0 \ar@{~>}[ld] \ar@{~>}[rd] &
& y'_1 \ar@{~>}[ld] & \ldots
& y'_{k - 1} \ar@{~>}[rd] &
\\
y_0 & &
y_1 & &
\ldots & &
y_k = y
}
$$
between points of $Y$ with $y_0 \in Z$ and $y_i' \leadsto y_i$
an immediate specialization, then $\delta_Z(y_k) \leq k$.
\end{enumerate}
\end{lemma}

\begin{proof}
Proof of (1). If $y \in Z$, then we can take $k = 0$ and
$V_0 = W_0 = \overline{\{y\}}$ and we get $\delta(V_0, W_0) = 0$ so
$\delta_Z(y) = 0$. If $y \not \in Z$, then for every system
$V_0 \subset W_0 \supset V_1 \subset W_1 \supset \ldots \subset W_k$
for $y$ we either have $k = 0$ and $V_0 \not = W_0$ or $k > 0$.
In both cases $\delta(V_0, W_0, \ldots, W_k) > 0$. Hence $\delta_Z(y) > 0$.

\medskip\noindent
Proof of (2). Let $y \leadsto y'$ be a nontrivial specialization and let
$V_0 \subset W_0 \supset V_1 \subset W_1 \supset \ldots \subset W_k$
is a system for $y$. Here there are two cases.
Case I: $V_k = W_k$, i.e., $c_k = 0$. In this case
we can set $V'_k = W'_k = \overline{\{y'\}}$.
An easy computation shows that
$\delta(V_0, W_0, \ldots, V'_k, W'_k) \leq
\delta(V_0, W_0, \ldots, V_k, W_k)$
because only $b_{k - 1}$ is changed into a bigger integer.
Case II: $V_k \not = W_k$, i.e., $c_k > 0$.
In this case, setting $V_{k + 1} = W_{k + 1} = \overline{\{y'\}}$
we see that $V_0 \subset W_0 \supset \ldots \subset W_k \supset V_{k + 1}
\subset W_{k + 1}$ is a system for $y'$. Then $c_{k + 1} = 0$ and
$b_k > 0$ so we get
\begin{align*}
& \delta(V_0, \ldots, W_{k + 1}) \\
& =
k + 1 +
\max_{i = 0, 1, \ldots, k + 1}
(c_i + c_{i + 1} + \ldots + c_k + c_{k + 1}
- b_i - b_{i + 1} - \ldots - b_{k - 1} - b_k) \\
& =
k + 1 +
\max_{i = 0, 1, \ldots, k + 1}
(c_i + c_{i + 1} + \ldots + c_k
- b_i - b_{i + 1} - \ldots - b_{k - 1} - b_k) \\
& \leq
k + \max_{i = 0, 1, \ldots, k}
(c_i + c_{i + 1} + \ldots + c_k - b_i - b_{i + 1} - \ldots - b_{k - 1}) \\
& = \delta(V_0, \ldots, W_k)
\end{align*}
The inequality holds because $c_k > 0$ and $b_k > 0$ which in particular
implies that $\delta(V_0, \ldots, W_k) \geq k + c_k \geq k + 1$.

\medskip\noindent
Proof of (3). Given $y \in Y$ and $z \in \overline{\{y\}} \cap Z$
we get the system
$$
V_0 = \overline{\{z\}} \subset W_0 = \overline{\{y\}}
$$
and $c_0 = \text{codim}(V_0, W_0) = \dim(\mathcal{O}_{\overline{\{y\}}, z})$
by Properties, Lemma \ref{properties-lemma-codimension-local-ring}.
Thus we see that $\delta(V_0, W_0) = 0 + c_0 = c_0$ which proves
what we want.

\medskip\noindent
Proof of (4). Let $\delta$ be a dimension function on $Y$.
Let $V_0 \subset W_0 \supset V_1 \subset W_1 \supset \ldots \subset W_k$
be a system for $y$. Let $y'_i \in W_i$ and $y_i \in V_i$ be the
generic points, so $y_0 \in Z$ and $y_k = y$. Then we see that
$$
\delta(y_i) - \delta(y_{i - 1}) =
\delta(y'_{i - 1}) - \delta(y_{i - 1}) - \delta(y'_{i - 1}) + \delta(y_i) =
c_{i - 1} - b_{i - 1}
$$
Finally, we have $\delta(y'_k) - \delta(y_{k - 1}) = c_k$.
Thus we see that
$$
\delta(y) - \delta(y_0) =
c_0 + \ldots + c_k - b_0 - \ldots - b_{k - 1}
$$
We conclude
$\delta(V_0, W_0, \ldots, W_k) \geq k + \delta(y) - \delta(y_0)$
which proves what we want.

\medskip\noindent
Proof of (5). The function $\delta(y) = \dim(\overline{\{y\}})$
is a dimension function. Hence $\delta(y) \leq \delta_Z(y)$ by
part (4). By part (3) we have $\delta_Z(y) \leq \delta(y)$
and we are done.

\medskip\noindent
Proof of (6). This is clear as their are fewer systems to consider
in the computation of $\delta^{Y'}_{Y' \cap Z}$.

\medskip\noindent
Proof of (7). Let
$V_0 \subset W_0 \supset V_1 \subset W_1 \supset \ldots \subset W_k$
be a system for $y$. Set $W'_k = \overline{\{y'\}}$.
Since $Y$ is catenary, we see that
$\text{codim}(V_k, W'_k) = \text{codim}(V_k, W_k) + 1$.
It follows easily that $\delta(V_0, \ldots, V_k, W'_k) \leq
\delta(V_0, \ldots, V_k, W_k) + 1$ which proves what we want.

\medskip\noindent
Proof of (8): combine (7) and (2).
\end{proof}

\begin{lemma}
\label{lemma-change-distance-function}
Let $Y$ be a universally catenary Noetherian scheme. Let $Z \subset Y$
be a closed subscheme. Let $f : Y' \to Y$ be a finite type
morphism all of whose fibres have dimension $\leq e$. Set $Z' = f^{-1}(Z)$.
Then
$$
\delta_Z(y) \leq \delta_{Z'}(y') + e - \text{trdeg}_{\kappa(y)}(\kappa(y'))
$$
for $y' \in Y'$ with image $y \in Y$.
\end{lemma}

\begin{proof}
If $\delta_{Z'}(y') = \infty$, then there is nothing to prove.
If $\delta_{Z'}(y') < \infty$, then we choose a system
of integral closed subschemes
$$
V'_0 \subset W'_0 \supset V'_1 \subset W'_1 \supset \ldots \subset W'_k
$$
of $Y'$ with $V'_0 \subset Z'$ and $y'$ the generic point of $W'_k$
such that $\delta_{Z'}(y') = \delta(V'_0, W'_0, \ldots, W'_k)$.
Denote
$$
V_0 \subset W_0 \supset V_1 \subset W_1 \supset \ldots \subset W_k
$$
the scheme theoretic images of the above schemes in $Y$. Observe
that $y$ is the generic point of $W_k$ and that $V_0 \subset Z$.
For each $i$ we look at the diagram
$$
\xymatrix{
V'_i \ar[r] \ar[d] & W'_i \ar[d] & V'_{i + 1} \ar[l] \ar[d] \\
V_i \ar[r] & W_i & V_{i + 1} \ar[l]
}
$$
Denote $n_i$ the relative dimension of $V'_i/V_i$ and
$m_i$ the relative dimension of $W'_i/W_i$; more precisely
these are the transcendence degrees of the corresponding extensions of
the function fields. Set
$c_i = \text{codim}(V_i, W_i)$,
$c'_i = \text{codim}(V'_i, W'_i)$,
$b_i = \text{codim}(V_{i + 1}, W_i)$, and
$b'_i = \text{codim}(V'_{i + 1}, W'_i)$.
By the dimension formula we have
$$
c_i = c'_i + n_i - m_i
\quad\text{and}\quad
b_i = b'_i + n_{i + 1} - m_i
$$
See Morphisms, Lemma \ref{morphisms-lemma-dimension-formula}.
Hence $c_i - b_i = c'_i - b'_i + n_i - n_{i + 1}$. Thus we see that
\begin{align*}
& c_i + c_{i + 1} + \ldots + c_k - b_i - b_{i + 1} - \ldots - b_{k - 1} \\
& =
c'_i + c'_{i + 1} + \ldots + c'_k - b'_i - b'_{i + 1} - \ldots - b'_{k - 1}
+ n_i - n_k + c_k - c'_k \\
& =
c'_i + c'_{i + 1} + \ldots + c'_k - b'_i - b'_{i + 1} - \ldots - b'_{k - 1}
+ n_i - m_k
\end{align*}
Thus we see that
\begin{align*}
\max_{i = 0, \ldots, k}
& (c_i + c_{i + 1} + \ldots + c_k - b_i - b_{i + 1} - \ldots - b_{k - 1}) \\
& =
\max_{i = 0, \ldots, k}
(c'_i + c'_{i + 1} + \ldots + c'_k - b'_i - b'_{i + 1} - \ldots - b'_{k - 1}
+ n_i - m_k) \\
& =
\max_{i = 0, \ldots, k}
(c'_i + c'_{i + 1} + \ldots + c'_k - b'_i - b'_{i + 1} - \ldots - b'_{k - 1}
+ n_i) - m_k \\
& \leq
\max_{i = 0, \ldots, k}
(c'_i + c'_{i + 1} + \ldots + c'_k - b'_i - b'_{i + 1} - \ldots - b'_{k - 1})
+ e - m_k
\end{align*}
Since $m_k = \text{trdeg}_{\kappa(y)}(\kappa(y'))$ we conclude that
$$
\delta(V_0, W_0, \ldots, W_k) \leq
\delta(V'_0, W'_0, \ldots, W'_k) + e - \text{trdeg}_{\kappa(y)}(\kappa(y'))
$$
as desired.
\end{proof}

\begin{remark}
\label{remark-discussion}
Let $Y$ be a catenary Noetherian scheme and let $Z \subset Y$
be a closed subset. By Lemma \ref{lemma-discussion} we have
$$
\delta_Z(y) \leq \min
\left\{ k \middle|
\begin{matrix}
\text{ there exist specializations in }Y \\
y_0 \leftarrow y'_0 \rightarrow y_1 \leftarrow y'_1 \rightarrow \ldots
\leftarrow y'_{k - 1} \rightarrow y_k = y \\
\text{ with }y_0 \in Z\text{ and }y_i' \leadsto y_i
\text{ immediate}
\end{matrix}
\right\}
$$
We claim that if $Y$ is of finite type over a field,
then equality holds. If we ever need this result we
will formulate a precise result and prove it here.
However, in general if we define $\delta_Z$
by the right hand side of this inequality, then we don't
know if Lemma \ref{lemma-change-distance-function} remains true.
\end{remark}

\begin{example}
\label{example-distance}
Let $k$ be a field and $Y = \mathbf{A}^n_k$. Denote
$\delta : Y \to \mathbf{Z}_{\geq 0}$ the usual dimension function.
\begin{enumerate}
\item If $Z = \{z\}$ for some closed point $z$, then
\begin{enumerate}
\item $\delta_Z(y) = \delta(y)$ if $y \leadsto z$ and
\item $\delta_Z(y) = \delta(y) + 1$ if $y \not \leadsto z$.
\end{enumerate}
\item If $Z$ is a closed subvariety and $W = \overline{\{y\}}$, then
\begin{enumerate}
\item $\delta_Z(y) = 0$ if $W \subset Z$,
\item $\delta_Z(y) = \dim(W) - \dim(Z)$ if $Z$ is contained in $W$,
\item $\delta_Z(y) = 1$ if $\dim(W) \leq \dim(Z)$ and $W \not \subset Z$,
\item $\delta_Z(y) = \dim(W) - \dim(Z) + 1$ if $\dim(W) > \dim(Z)$
and $Z \not \subset W$.
\end{enumerate}
\end{enumerate}
A generalization of case (1) is if $Y$ is of finite type over a field
and $Z = \{z\}$ is a closed point. Then $\delta_Z(y) = \delta(y) + t$
where $t$ is the minimum length of a chain of curves connecting
$z$ to a closed point of $\overline{\{y\}}$.
\end{example}







\section{Formal sections, III}
\label{section-formal-sections-cd-one}

\noindent
In this section we prove Lemma \ref{lemma-topology-I-adic}
which (in the setting of Noetherian schemes and coherent modules)
is the analogue of Cohomology, Lemma \ref{cohomology-lemma-topology-I-adic-f}
in case the ideal $I$ is not assumed principal but
has the property that $\text{cd}(A, I) = 1$.

\begin{lemma}
\label{lemma-cd-one}
Let $I = (f_1, \ldots, f_r)$ be an ideal of a Noetherian ring $A$.
If $\text{cd}(A, I) = 1$, then there exist $c \geq 1$ and maps
$\varphi_j : I^c \to A$ such that $\sum f_j \varphi_j : I^c \to I$
is the inclusion map.
\end{lemma}

\begin{proof}
Since $\text{cd}(A, I) = 1$ the complement $U = \Spec(A) \setminus V(I)$
is affine (Local Cohomology, Lemma \ref{local-cohomology-lemma-cd-is-one}).
Say $U = \Spec(B)$. Then $IB = B$
and we can write $1 = \sum_{j = 1, \ldots, r} f_j b_j$
for some $b_j \in B$. By
Cohomology of Schemes, Lemma \ref{coherent-lemma-homs-over-open}
we can represent $b_j$ by maps $\varphi_j : I^c \to A$
for some $c \geq 0$. Then $\sum f_j \varphi_j : I^c \to I \subset A$
is the canonical embedding, after possibly replacing $c$ by a larger
integer, by the same lemma.
\end{proof}

\begin{lemma}
\label{lemma-cd-one-extend}
Let $I = (f_1, \ldots, f_r)$ be an ideal of a Noetherian ring $A$
with $\text{cd}(A, I) = 1$. Let $c \geq 1$ and $\varphi_j : I^c \to A$,
$j = 1, \ldots, r$ be as in Lemma \ref{lemma-cd-one}.
Then there is a unique graded $A$-algebra map
$$
\Phi : \bigoplus\nolimits_{n \geq 0} I^{nc} \to A[T_1, \ldots, T_r]
$$
with $\Phi(g) = \sum \varphi_j(g) T_j$ for $g \in I^c$.
Moreover, the composition of $\Phi$ with the map
$A[T_1, \ldots, T_r] \to \bigoplus_{n \geq 0} I^n$,
$T_j \mapsto f_j$ is the inclusion map
$\bigoplus_{n \geq 0} I^{nc} \to \bigoplus_{n \geq 0} I^n$.
\end{lemma}

\begin{proof}
For each $j$ and $m \geq c$ the restriction of $\varphi_j$ to
$I^m$ is a map $\varphi_j : I^m \to I^{m - c}$.
Given $j_1, \ldots, j_n \in \{1, \ldots, r\}$ we claim that the
composition
$$
\varphi_{j_1} \ldots \varphi_{j_n} :
I^{nc} \to I^{(n - 1)c} \to \ldots \to I^c \to A
$$
is independent of the order of the indices $j_1, \ldots, j_n$.
Namely, if $g = g_1 \ldots g_n$ with $g_i \in I^c$, then
we see that
$$
(\varphi_{j_1} \ldots \varphi_{j_n})(g) =
\varphi_{j_1}(g_1) \ldots \varphi_{j_n}(g_n)
$$
is independent of the ordering as multiplication in $A$ is commutative.
Thus we can define $\Phi$ by sending $g \in I^{nc}$ to
$$
\Phi(g) = \sum\nolimits_{e_1 + \ldots + e_r = n}
(\varphi_1^{e_1} \circ \ldots \circ \varphi_r^{e_r})(g)
T_1^{e_1} \ldots T_r^{e_r}
$$
It is straightforward to prove that this is a graded $A$-algebra
homomorphism with the desired property. Uniqueness is immediate
as is the final property. This proves the lemma.
\end{proof}

\begin{lemma}
\label{lemma-cd-one-extend-to-module}
Let $I = (f_1, \ldots, f_r)$ be an ideal of a Noetherian ring $A$
with $\text{cd}(A, I) = 1$. Let $c \geq 1$ and $\varphi_j : I^c \to A$,
$j = 1, \ldots, r$ be as in Lemma \ref{lemma-cd-one}.
Let $A \to B$ be a ring map with $B$ Noetherian and let $N$ be
a finite $B$-module. Then, after possibly increasing $c$
and adjusting $\varphi_j$ accordingly, there is a unique
unique graded $B$-module map
$$
\Phi_N : \bigoplus\nolimits_{n \geq 0} I^{nc}N \to N[T_1, \ldots, T_r]
$$
with $\Phi_N(g x) = \Phi(g) x$ for $g \in I^{nc}$ and $x \in N$
where $\Phi$ is as in Lemma \ref{lemma-cd-one-extend}.
The composition of $\Phi_N$ with the map
$N[T_1, \ldots, T_r] \to \bigoplus_{n \geq 0} I^nN$,
$T_j \mapsto f_j$ is the inclusion map
$\bigoplus_{n \geq 0} I^{nc}N \to \bigoplus_{n \geq 0} I^nN$.
\end{lemma}

\begin{proof}
The uniqueness is clear from the formula and the uniqueness of $\Phi$ in
Lemma \ref{lemma-cd-one-extend}. Consider the Noetherian $A$-algebra
$B' = B \oplus N$ where $N$ is an ideal of square zero. To show
the existence of $\Phi_N$ it is enough
(via Lemma \ref{lemma-cd-one}) to show that $\varphi_j$ extends to
a map $\varphi'_j : I^cB' \to B'$ after possibly increasing $c$
to some $c'$ (and replacing $\varphi_j$ by the composition of the inclusion
$I^{c'} \to I^c$ with $\varphi_j$). Recall that $\varphi_j$ corresponds to a
section
$$
h_j \in \Gamma(\Spec(A) \setminus V(I), \mathcal{O}_{\Spec(A)})
$$
see Cohomology of Schemes, Lemma \ref{coherent-lemma-homs-over-open}.
(This is in fact how we chose our $\varphi_j$ in the proof of
Lemma \ref{lemma-cd-one}.) Let us use the same lemma to represent the pullback
$$
h'_j \in \Gamma(\Spec(B') \setminus V(IB'), \mathcal{O}_{\Spec(B')})
$$
of $h_j$ by a $B'$-linear map
$\varphi'_j : I^{c'}B' \to B'$ for some $c' \geq c$.
The agreement with $\varphi_j$ will hold for $c'$
sufficiently large by a further application of the lemma:
namely we can test agreement on a finite list of generators of $I^{c'}$.
Small detail omitted.
\end{proof}

\begin{lemma}
\label{lemma-cd-is-one-for-system}
Let $I = (f_1, \ldots, f_r)$ be an ideal of a Noetherian ring $A$ with
$\text{cd}(A, I) = 1$. Let $c \geq 1$ and $\varphi_j : I^c \to A$,
$j = 1, \ldots, r$ be as in Lemma \ref{lemma-cd-one}.
Let $X$ be a Noetherian scheme over $\Spec(A)$. Let
$$
\ldots \to \mathcal{F}_3 \to \mathcal{F}_2 \to \mathcal{F}_1
$$
be an inverse system of coherent $\mathcal{O}_X$-modules
such that $\mathcal{F}_n = \mathcal{F}_{n + 1}/I^n\mathcal{F}_{n + 1}$.
Set $\mathcal{F} = \lim \mathcal{F}_n$.
Then, after possibly increasing $c$ and adjusting $\varphi_j$ accordingly,
there exists a unique graded $\mathcal{O}_X$-module map
$$
\Phi_\mathcal{F} :
\bigoplus\nolimits_{n \geq 0} I^{nc}\mathcal{F}
\longrightarrow
\mathcal{F}[T_1, \ldots, T_r]
$$
with $\Phi_\mathcal{F}(g s) = \Phi(g) s$ for $g \in I^{nc}$ and
$s$ a local section of $\mathcal{F}$ where $\Phi$ is as in
Lemma \ref{lemma-cd-one-extend}. The composition of $\Phi_\mathcal{F}$
with the map
$\mathcal{F}[T_1, \ldots, T_r] \to \bigoplus_{n \geq 0} I^n\mathcal{F}$,
$T_j \mapsto f_j$
is the canonical inclusion
$\bigoplus_{n \geq 0} I^{nc}\mathcal{F} \to
\bigoplus_{n \geq 0} I^n\mathcal{F}$.
\end{lemma}

\begin{proof}
The uniqueness is immediate from the $\mathcal{O}_X$-linearity
and the requirement that $\Phi_\mathcal{F}(g s) = \Phi(g) s$ for
$g \in I^{nc}$ and $s$ a local section of $\mathcal{F}$.
Thus we may assume $X = \Spec(B)$ is affine.
Observe that $(\mathcal{F}_n)$ is an object of the category
$\textit{Coh}(X, I\mathcal{O}_X)$ introduced
in Cohomology of Schemes, Section \ref{coherent-section-coherent-formal}.
Let $B' = B^\wedge$ be the $I$-adic completion of $B$.
By Cohomology of Schemes, Lemma \ref{coherent-lemma-inverse-systems-affine}
the object $(\mathcal{F}_n)$ corresponds to a finite $B'$-module $N$
in the sense that $\mathcal{F}_n$ is the coherent
module associated to the finite $B$-module $N/I^n N$.
Applying Lemma \ref{lemma-cd-one-extend-to-module}
to $I \subset A \to B'$ and $N$
we see that, after possibly increasing $c$ and adjusting
$\varphi_j$ accordingly, we get unique maps
$$
\Phi_N : \bigoplus\nolimits_{n \geq 0} I^{nc}N \to N[T_1, \ldots, T_r]
$$
with the corresponding properties. Note that in degree $n$ we obtain
an inverse system of maps $N/I^mN \to \bigoplus_{e_1 + \ldots + e_r = n}
N/I^{m - nc}N \cdot T_1^{e_1} \ldots T_r^{e_r}$ for $m \geq nc$.
Translating back into coherent
sheaves we see that $\Phi_N$ corresponds to a system of maps
$$
\Phi^n_m :
I^{nc}\mathcal{F}_m
\longrightarrow
\bigoplus\nolimits_{e_1 + \ldots + e_r = n}
\mathcal{F}_{m - nc} \cdot T_1^{e_1} \ldots T_r^{e_r}
$$
for varying $m \geq nc$ and $n \geq 1$. Taking the inverse limit of
these maps over $m$ we obtain $\Phi_\mathcal{F} = \bigoplus_n \lim_m \Phi^n_m$.
Note that $\lim_m I^t\mathcal{F}_m = I^t \mathcal{F}$ as can be seen by
evaluating on affines for example, but in fact we don't need this because
it is clear there is a map $I^t\mathcal{F} \to \lim_m I^t\mathcal{F}_m$.
\end{proof}

\begin{lemma}
\label{lemma-topology-I-adic}
Let $I$ be an ideal of a Noetherian ring $A$. Let $X$ be a Noetherian scheme
over $\Spec(A)$. Let
$$
\ldots \to \mathcal{F}_3 \to \mathcal{F}_2 \to \mathcal{F}_1
$$
be an inverse system of coherent $\mathcal{O}_X$-modules
such that $\mathcal{F}_n = \mathcal{F}_{n + 1}/I^n\mathcal{F}_{n + 1}$.
If $\text{cd}(A, I) = 1$, then for all $p \in \mathbf{Z}$ the limit topology on
$\lim H^p(X, \mathcal{F}_n)$ is $I$-adic.
\end{lemma}

\begin{proof}
First it is clear that $I^t \lim H^p(X, \mathcal{F}_n)$
maps to zero in $H^p(X, \mathcal{F}_t)$. Thus the $I$-adic topology
is finer than the limit topology. For the converse we set
$\mathcal{F} = \lim \mathcal{F}_n$, we pick generators $f_1, \ldots, f_r$
of $I$, we pick $c \geq 1$, and we choose
$\Phi_\mathcal{F}$ as in Lemma \ref{lemma-cd-is-one-for-system}.
We will use the results of Lemma \ref{lemma-properties-system}
without further mention. In particular we have a short exact
sequence
$$
0 \to R^1\lim H^{p - 1}(X, \mathcal{F}_n) \to H^p(X, \mathcal{F})
\to \lim H^p(X, \mathcal{F}_n) \to 0
$$
Thus we can lift any element $\xi$ of $\lim H^p(X, \mathcal{F}_n)$
to an element $\xi' \in H^p(X, \mathcal{F})$. Suppose $\xi$ maps to zero
in $H^p(X, \mathcal{F}_{nc})$ for some $n$, in other
words, suppose $\xi$ is ``small'' in the limit topology. We have a
short exact sequence
$$
0 \to I^{nc}\mathcal{F} \to \mathcal{F} \to \mathcal{F}_{nc} \to 0
$$
and hence the assumption means we can lift $\xi'$ to an element
$\xi'' \in H^p(X, I^{nc}\mathcal{F})$. Applying $\Phi_\mathcal{F}$
we get
$$
\Phi_\mathcal{F}(\xi'') = \sum\nolimits_{e_1 + \ldots + e_r = n}
\xi'_{e_1, \ldots, e_r} \cdot T_1^{e_1} \ldots T_r^{e_r}
$$
for some $\xi'_{e_1, \ldots, e_r} \in H^p(X, \mathcal{F})$.
Letting $\xi_{e_1, \ldots, e_r} \in \lim H^p(X, \mathcal{F}_n)$
be the images and using the final assertion of
Lemma \ref{lemma-cd-is-one-for-system}
we conclude that
$$
\xi = \sum f_1^{e_1} \ldots f_r^{e_r} \xi_{e_1, \ldots, e_r}
$$
is in $I^n \lim H^p(X, \mathcal{F}_n)$ as desired.
\end{proof}

\begin{example}
\label{example-not-I-adic}
Let $k$ be a field. Let $A = k[x, y][[s, t]]/(xs - yt)$.
Let $I = (s, t)$ and $\mathfrak a = (x, y, s, t)$.
Let $X = \Spec(A) - V(\mathfrak a)$ and
$\mathcal{F}_n = \mathcal{O}_X/I^n\mathcal{O}_X$.
Observe that the rational function
$$
g = \frac{t}{x} = \frac{s}{y}
$$
is regular in an open neighbourhood $V \subset X$ of
$V(I\mathcal{O}_X)$. Hence every power $g^e$ determines a section
$g^e \in M = \lim H^0(X, \mathcal{F}_n)$. Observe that
$g^e \to 0$ as $e \to \infty$ in the limit topology on $M$
since $g^e$ maps to zero in $\mathcal{F}_e$.
On the other hand, $g^e \not \in IM$ for any $e$
as the reader can see by computing $H^0(U, \mathcal{F}_n)$;
computation omitted. Observe that $\text{cd}(A, I) = 2$.
Thus the result of Lemma \ref{lemma-topology-I-adic} is sharp.
\end{example}







\begin{lemma}
\label{lemma-when-done}
In Situation \ref{situation-algebraize} let $(\mathcal{F}_n)$ be an object
of $\textit{Coh}(U, I\mathcal{O}_U)$. Let $M$ be as in (\ref{equation-guess}).
Assume
\begin{enumerate}
\item[(a)] the inverse system $H^0(U, \mathcal{F}_n)$ has Mittag-Leffler,
\item[(b)] the limit topology on $M$ agrees with the $I$-adic topology, and
\item[(c)] the image of $M \to H^0(U, \mathcal{F}_n)$ is a finite $A$-module
for all $n$.
\end{enumerate}
Then $(\mathcal{F}_n)$ extends canonically to $X$.
In particular, if $A$ is $I$-adically complete, then
$(\mathcal{F}_n)$ is the completion of a coherent $\mathcal{O}_U$-module.
\end{lemma}

\begin{proof}
Since $H^0(U, \mathcal{F}_n)$ has the Mittag-Leffler condition
and since the limit topology on $M$ is the $I$-adic topology
we see that $\{M/I^nM\}$ and $\{H^0(U, \mathcal{F}_n)\}$
are pro-isomorphic inverse systems of $A$-modules.
Thus if we set
$$
\mathcal{G}_n = \widetilde{M/I^n M}
$$
then we see that to verify the condition in
Definition \ref{definition-canonically-algebraizable}
it suffices to show that $M$ is a finite module over the
$I$-adic completion of $A$. This follows from
the fact that $M/I^n M$ is finite by condition (c)
and the above and
Algebra, Lemma \ref{algebra-lemma-finite-over-complete-ring}.
\end{proof}

\noindent
The following is in some sense the most straightforward possible application
Lemma \ref{lemma-when-done} above.

\begin{lemma}
\label{lemma-algebraization-principal-variant}
In Situation \ref{situation-algebraize} let
$(\mathcal{F}_n)$ be an object of $\textit{Coh}(U, I\mathcal{O}_U)$.
Assume
\begin{enumerate}
\item $I = (f)$ is a principal ideal for a nonzerodivisor $f \in \mathfrak a$,
\item $\mathcal{F}_n$ is a finite locally free
$\mathcal{O}_U/f^n\mathcal{O}_U$-module,
\item $H^1_\mathfrak a(A/fA)$ and $H^2_\mathfrak a(A/fA)$
are finite $A$-modules.
\end{enumerate}
Then $(\mathcal{F}_n)$ extends canonically to $X$. In particular, if $A$
is complete, then $(\mathcal{F}_n)$ is the completion of a coherent
$\mathcal{O}_U$-module.
\end{lemma}

\begin{proof}
We will prove this by verifying hypotheses (a), (b), and (c) of
Lemma \ref{lemma-when-done}.

\medskip\noindent
Since $\mathcal{F}_n$ is locally free over $\mathcal{O}_U/f^n\mathcal{O}_U$
we see that we have short exact sequences
$0 \to \mathcal{F}_n \to \mathcal{F}_{n + 1} \to \mathcal{F}_1 \to 0$
for all $n$. Thus condition (b) holds by
Cohomology, Lemma \ref{cohomology-lemma-topology-I-adic-f}.

\medskip\noindent
As $f$ is a nonzerodivisor we obtain short exact sequences
$$
0 \to A/f^nA \xrightarrow{f} A/f^{n + 1}A \to A/fA \to 0
$$
and we have corresponding short exact sequences
$0 \to \mathcal{F}_n \to \mathcal{F}_{n + 1} \to \mathcal{F}_1 \to 0$.
We will use
Local Cohomology, Lemma
\ref{local-cohomology-lemma-finiteness-pushforwards-and-H1-local}
without further mention. Our assumptions imply that
$H^0(U, \mathcal{O}_U/f\mathcal{O}_U)$ and
$H^1(U, \mathcal{O}_U/f\mathcal{O}_U)$
are finite $A$-modules. Hence the same thing is true for $\mathcal{F}_1$, see
Local Cohomology, Lemma
\ref{local-cohomology-lemma-finiteness-for-finite-locally-free}.
Using induction and the short exact sequences we find that
$H^0(U, \mathcal{F}_n)$ are finite $A$-modules for all $n$.
In this way we see hypothesis (c) is satisfied.

\medskip\noindent
Finally, as $H^1(U, \mathcal{F}_1)$ is a finite $A$-module
we can apply Cohomology, Lemma \ref{cohomology-lemma-ML}
to see hypothesis (a) holds.
\end{proof}

\begin{remark}
\label{remark-interesting-case-variant}
In Lemma \ref{lemma-algebraization-principal-variant}
if $A$ is universally catenary with Cohen-Macaulay
formal fibres (for example if $A$ has a dualizing complex), then
the condition that
$H^1_\mathfrak a(A/fA)$ and $H^2_\mathfrak a(A/fA)$
are finite $A$-modules, is equivalent with
$$
\text{depth}((A/f)_\mathfrak p) + \dim((A/\mathfrak p)_\mathfrak q) > 2
$$
for all $\mathfrak p \in V(f) \setminus V(\mathfrak a)$
and $\mathfrak q \in V(\mathfrak p) \cap V(\mathfrak a)$
by Local Cohomology, Theorem \ref{local-cohomology-theorem-finiteness}.

\medskip\noindent
For example, if $A/fA$ is $(S_2)$ and if every irreducible
component of $Z = V(\mathfrak a)$ has codimension $\geq 3$
in $Y = \Spec(A/fA)$, then we get the finiteness of
$H^1_\mathfrak a(A/fA)$ and $H^2_\mathfrak a(A/fA)$.
This should be contrasted with the slightly weaker conditions
found in Lemma \ref{lemma-algebraization-principal}
(see also Remark \ref{remark-interesting-case}).
\end{remark}






\begin{proof}
We will check hypotheses (a), (b), and (c) of Lemma \ref{lemma-when-done}.
Before we start, let us point out that the modules
$H^0(U, \mathcal{F}''_n)$ and $H^1(U, \mathcal{F}''_n)$
are finite $A$-modules for all $n$ by
Local Cohomology, Lemma \ref{local-cohomology-lemma-finiteness-Rjstar}.

\medskip\noindent
Observe that for each $p \geq 0$
the limit topology on $\lim H^p(U, \mathcal{F}_n)$
is the $I$-adic topology by Lemma \ref{lemma-topology-I-adic}.
In particular, hypothesis (b) holds.

\medskip\noindent
We know that $M = \lim H^0(U, \mathcal{F}_n)$ is an $A$-module whose
limit topology is the $I$-adic topology. Thus, given $n$, the module
$M/I^nM$ is a subquotient of $H^0(U, \mathcal{F}_N)$ for some $N \gg n$.
Since the inverse system $\{H^0(U, \mathcal{F}_N)\}$ is pro-isomorphic to an
inverse system of finite $A$-modules, namely $\{H^0(U, \mathcal{F}''_N)\}$,
we conclude that $M/I^nM$ is finite. It follows that $M$ is finite, see
Algebra, Lemma \ref{algebra-lemma-finite-over-complete-ring}.
In particular hypothesis (c) holds.

\medskip\noindent
For each $n \geq 0$ let us write $Ob_n = \lim_N H^1(U, I^n\mathcal{F}_N)$.
A special case is $Ob = Ob_0 = \lim_N H^1(U, \mathcal{F}_N)$.
Arguing exactly as in the previous paragraph we find that $Ob$
is a finite $A$-module. (In fact, we also know that $Ob/I Ob$ is annihilated
by a power of $\mathfrak a$, but it seems somewhat difficult to use this.)

\medskip\noindent
We set $\mathcal{F} = \lim \mathcal{F}_n$, we pick generators
$f_1, \ldots, f_r$ of $I$, we pick $c \geq 1$, and we choose
$\Phi_\mathcal{F}$ as in Lemma \ref{lemma-cd-is-one-for-system}.
We will use the results of Lemma \ref{lemma-properties-system}
without further mention. In particular, for each $n \geq 1$ there are maps
$$
\delta_n :
H^0(U, \mathcal{F}_n)
\longrightarrow
H^1(U, I^n\mathcal{F})
\longrightarrow
Ob_n
$$
The first comes from the short exact sequence
$0 \to I^n\mathcal{F} \to \mathcal{F} \to \mathcal{F}_n \to 0$
and the second from $I^n\mathcal{F} = \lim I^n\mathcal{F}_N$.
We will later use that if $\delta_n(s) = 0$ for $s \in H^0(U, \mathcal{F}_n)$
then we can for each $n' \geq n$ find $s' \in H^0(U, \mathcal{F}_{n'})$
mapping to $s$.
Observe that there are commutative diagrams
$$
\xymatrix{
H^0(U, \mathcal{F}_{nc}) \ar[r] \ar[dd] &
H^1(U, I^{nc}\mathcal{F}) \ar[dd] \ar[rd]^{\Phi_\mathcal{F}} \\
& &
\bigoplus_{e_1 + \ldots + e_r = n}
H^1(U, \mathcal{F}) \cdot T_1^{e_1} \ldots T_r^{e_r} \ar[ld] \\
H^0(U, \mathcal{F}_n) \ar[r] &
H^1(U, I^n\mathcal{F})
}
$$
We conclude that the obstruction map
$H^0(U, \mathcal{F}_n) \to Ob_n$
sends the image of
$H^0(U, \mathcal{F}_{nc}) \to H^0(U, \mathcal{F}_n)$
into the submodule
$$
Ob'_n =
\Im\left(
\bigoplus\nolimits_{e_1 + \ldots + e_r = n}
Ob \cdot T_1^{e_1} \ldots T_r^{e_r} \to Ob_n
\right)
$$
where on the summand $Ob \cdot T_1^{e_1} \ldots T_r^{e_r}$
we use the map on cohomology coming from the reductions modulo
powers of $I$ of the multiplication map
$f_1^{e_1} \ldots f_r^{e_r} : \mathcal{F} \to I^n\mathcal{F}$.
By construction
$$
\bigoplus\nolimits_{n \geq 0} Ob'_n
$$
is a finite graded module over the Rees algebra $\bigoplus_{n \geq 0} I^n$.
For each $n$ we set
$$
M_n = \{s \in H^0(U, \mathcal{F}_n) \mid \delta_n(s) \in Ob'_n\}
$$
Observe that $\{M_n\}$ is an inverse system and that
$f_j : \mathcal{F}_n \to \mathcal{F}_{n + 1}$ on global
sections maps $M_n$ into $M_{n + 1}$.
By exactly the same argument as in the proof of
Cohomology, Lemma \ref{cohomology-lemma-ML-general}
we find that $\{M_n\}$ is ML. Namely, because the Rees algebra
is Noetherian we can choose a finite number of homogeneous generators
of the form $\delta_{n_j}(z_j)$ with $z_j \in M_{n_j}$ for the graded submodule
$\bigoplus_{n \geq 0} \Im(M_n \to Ob'_n)$.
Then if $k = \max(n_j)$ we find that for $n \geq k$
and any $z \in M_n$ we can find $a_j \in I^{n - n_j}$ such that
$z - \sum a_j z_j$ is in the kernel of $\delta_n$
and hence in the image of $M_{n'}$ for all $n' \geq n$
(because the vanishing of $\delta_n$ means that we can
lift $z - \sum a_j z_j$ to an element $z' \in H^0(U, \mathcal{F}_{n'c})$
for all $n' \ge n$ and then the image of $z'$ in $H^0(U, \mathcal{F}_{n'})$
is in $M_{n'}$ by what we proved above).
Thus $\Im(M_n \to M_{n - k}) = \Im(M_{n'} \to M_{n - k})$
for all $n' \geq n$.

\medskip\noindent
Choose $n$. By the Mittag-Leffler property of $\{M_n\}$ we just established
we can find an $n' \geq n$ such that the image of $M_{n'} \to M_n$
is the same as the image of $M' \to M_n$. By the above we see that
the image of $M' \to M_n$ contains the image of
$H^0(U, \mathcal{F}_{n'c}) \to H^0(U, \mathcal{F}_n)$.
Thus we see that $\{M_n\}$ and $\{H^0(U, \mathcal{F}_n)\}$
are pro-isomorphic. Therefore $\{H^0(U, \mathcal{F}_n)\}$
has ML and we finally conclude that hypothesis (a) holds.
This concludes the proof.
\end{proof}

\bibliographystyle{plain}
\bibliography{my}
\end{document}
